\documentclass[a4paper,12pt]{article}
\usepackage[T2A]{fontenc}
\usepackage[utf8]{inputenc}
\usepackage[russian]{babel}
\usepackage{amsmath,amssymb,mathtools}
\usepackage{helvet}
\renewcommand{\familydefault}{\sfdefault}
\usepackage{icomma}
\usepackage{geometry}
\geometry{left=18mm,right=18mm,top=18mm,bottom=19mm}
\usepackage{microtype}
\usepackage{tikz}
\usetikzlibrary{calc,arrows.meta,positioning,shapes.geometric}
\usepackage{tabularx,booktabs,enumitem}
\usepackage{hyperref}
\usepackage{parskip}
\usepackage{fancyhdr}
\usepackage{setspace}
\usepackage{xcolor}

\definecolor{accent}{HTML}{2B6F73}
\definecolor{darkaccent}{HTML}{19484B}
\definecolor{textgray}{HTML}{2E3335}
\definecolor{softgray}{HTML}{6D777A}
\definecolor{linegray}{HTML}{D8DEDF}
\hypersetup{colorlinks=true,linkcolor=darkaccent,urlcolor=darkaccent}
\onehalfspacing
\setlength{\parindent}{0pt}
\setlist[itemize]{leftmargin=6mm,itemsep=2pt,topsep=3pt}
\setlist[enumerate]{leftmargin=8mm,itemsep=4pt,topsep=4pt}
\setlength{\headheight}{14.5pt}
\pagestyle{fancy}
\fancyhf{}
\fancyhead[L]{\small\color{softgray}София Хоменко}
\fancyhead[R]{\small\color{softgray}13.10.26}
\fancyfoot[C]{\small\color{softgray}\thepage}
\renewcommand{\headrulewidth}{0pt}
\renewcommand{\footrulewidth}{0pt}

\newcommand{\topic}[1]{%
  \vspace{0.45em}\noindent{\Large\bfseries\color{darkaccent}#1}\par
  \vspace{0.2em}\noindent\color{linegray}\rule{\linewidth}{0.6pt}\color{textgray}\vspace{0.3em}
}
\newcommand{\key}[1]{\textbf{\color{darkaccent}#1}}
\newcommand{\warn}{\textbf{\color{darkaccent}Важно: }}
\newcommand{\solution}[1]{\medskip\noindent\textbf{Пример.} #1}

\begin{document}
\color{textgray}
\thispagestyle{empty}
\begin{center}
\vspace*{24mm}
{\Large\color{softgray}Математика \textbullet\ ОГЭ}\par
\vspace{10mm}
{\Huge\bfseries\color{darkaccent}Чек-лист}\par
\vspace{7mm}
{\LARGE\bfseries Уравнения повышенной сложности}\par
\vspace{4mm}
{\Large Группировка, вынесение множителей,\par формулы сокращённого умножения}\par
\vspace{17mm}
{\large София Хоменко}\par
\vspace{4mm}
{\normalsize 13 октября 2026 года}\par
\vfill
{\small Лёвин Артём Александрович эксклюзивно для Софии Хоменко}\par
\vspace{12mm}
\begin{tikzpicture}[x=1cm,y=1cm]
  \draw[accent,line width=1.2pt]
  (0,0) .. controls (3.2,0.8) and (7.3,-0.8) .. (10.5,0)
        .. controls (13.1,0.6) and (15.6,0.2) .. (17.3,-0.15);
\end{tikzpicture}
\end{center}
\clearpage

\topic{1. Сначала приведите уравнение к нулю}

Основной инструмент сегодняшней темы --- \key{разложение на множители}. Если справа стоит выражение, перенесите все слагаемые влево \emph{со сменой знака}, чтобы получить
\[
F(x)=0.
\]

\key{Правило нулевого произведения:}
\[
A(x)\cdot B(x)=0
\quad\Longleftrightarrow\quad
A(x)=0\ \text{или}\ B(x)=0.
\]

\warn правило работает, только если \emph{вся левая часть} --- произведение, равное нулю. Для \((x-2)(x+3)+1=0\) приравнивать множители к нулю нельзя: сначала придётся преобразовать выражение.

\topic{2. Группировка и двойное вынесение}

Если слева четыре слагаемых, попробуйте объединить их в две группы. Цель --- получить \emph{одинаковый двучлен} в обеих группах. Само по себе наличие четырёх слагаемых ещё не гарантирует успех.

\solution{Решим уравнение}
\[
x^3-3x^2-8x+24=0.
\]
\begin{align*}
(x^3-3x^2)+(-8x+24)&=0,\\
x^2(x-3)-8(x-3)&=0,\\
(x-3)(x^2-8)&=0.
\end{align*}
Следовательно,
\[
x=3\qquad\text{или}\qquad x^2=8.
\]
Отсюда \(x=3\) либо \(x=\pm2\sqrt2\).

\key{Ответ:} \(-2\sqrt2;\ 2\sqrt2;\ 3\).

\warn при группировке отрицательное слагаемое должно сохранить знак:
\[
-8x+24=-8(x-3).
\]
\topic{3. Выносите общий множитель --- число, переменную, скобку}

Если каждый член содержит \(x\), вынесите его за скобки:
\[
x^3-4x^2-5x=x(x^2-4x-5)=x(x-5)(x+1).
\]
Уравнение \(x^3-4x^2-5x=0\) имеет корни \(\boxed{-1;\ 0;\ 5}\).

\warn нельзя сокращать обе части уравнения на \(x\), пока не проверен случай \(x=0\): так теряется корень.

\topic{4. Общим множителем бывает целое произведение}

\solution{Рассмотрим разность произведений:}
\begin{align*}
&(x-3)(x-4)(x-5)-(x-2)(x-4)(x-5)=0,\\[2pt]
&(x-4)(x-5)\bigl[(x-3)-(x-2)\bigr]=0,\\
&-(x-4)(x-5)=0.
\end{align*}
\key{Ответ:} \(4;\ 5\).

\warn минус перед скобкой меняет знак \emph{каждого} её слагаемого:
\[
(x-3)-(x-2)=x-3-x+2=-1.
\]

\clearpage
\topic{5. Узнавайте формулы сокращённого умножения}

Прежде чем раскрывать скобки, посмотрите, нельзя ли \emph{свернуть} выражение:
\[
a^2+2ab+b^2=(a+b)^2,\qquad
 a^2-2ab+b^2=(a-b)^2.
\]

\solution{Решим уравнение \(x(x^2+2x+1)=2(x+1)\).}
\begin{align*}
x(x+1)^2-2(x+1)&=0,\\
(x+1)\bigl[x(x+1)-2\bigr]&=0,\\
(x+1)(x^2+x-2)&=0,\\
(x+1)(x+2)(x-1)&=0.
\end{align*}
\key{Ответ:} \(-2;\ -1;\ 1\).

\warn не делите уравнение на \(x+1\): потеряете корень \(x=-1\). Общий множитель следует \emph{вынести}, затем применить правило нулевого произведения.

\topic{6. Разность квадратов: сохраняйте удобные скобки}

\[
\boxed{\displaystyle a^2-b^2=(a-b)(a+b).}
\]

\solution{Решим \(x^4=(4x-5)^2\). Здесь \(x^4=(x^2)^2\).}
\begin{align*}
(x^2)^2-(4x-5)^2&=0,\\
\bigl[x^2-(4x-5)\bigr]\bigl[x^2+(4x-5)\bigr]&=0,\\
(x^2-4x+5)(x^2+4x-5)&=0.
\end{align*}
Первое квадратное уравнение действительных корней не имеет:
\(D=16-20=-4\). Во втором \(D=16+20=36\), корни \(-5\) и \(1\).

\key{Ответ:} \(-5;\ 1\).

\clearpage
\topic{7. Разность квадратов экономит вычисления}

Если \(A(x)^2=B(x)^2\), удобно записать
\[
\bigl(A(x)-B(x)\bigr)\bigl(A(x)+B(x)\bigr)=0.
\]
При этом \emph{не нужно} раскрывать оба квадрата многочленов.

\solution{В уравнении \((2x-3)^2=(1-2x)^2\) получаем}
\begin{align*}
\bigl[(2x-3)-(1-2x)\bigr]
\bigl[(2x-3)+(1-2x)\bigr]&=0,\\
(4x-4)\cdot(-2)&=0.
\end{align*}
\key{Ответ:} \(1\).

\warn когда удалось получить произведение множителей, раскрывать его снова не требуется. Сразу решайте уравнения, в которых каждый множитель равен нулю.

\topic{8. Сумма квадратов равна нулю: нужна система}

Для любых действительных \(A\) и \(B\): \(A^2\ge0\), \(B^2\ge0\). Поэтому
\[
\boxed{\displaystyle A(x)^2+B(x)^2=0
\quad\Longleftrightarrow\quad
\begin{cases}A(x)=0,\\B(x)=0.\end{cases}}
\]

\solution{Решим уравнение}
\[
(x^2-25)^2+(x^2+3x-10)^2=0.
\]
Оба квадрата обязаны одновременно обратиться в нуль:
\[
\begin{cases}
x^2-25=0,\\x^2+3x-10=0.
\end{cases}
\quad\Longleftrightarrow\quad
\begin{cases}
x\in\{-5,5\},\\x\in\{-5,2\}.
\end{cases}
\]
Общим корнем служит только \(-5\).

\key{Ответ:} \(-5\).

\warn здесь используется \emph{система} (условия выполняются вместе); при нулевом \emph{произведении} --- совокупность (достаточно хотя бы одного нулевого множителя).

\clearpage
\thispagestyle{fancy}
\begin{center}
{\LARGE\bfseries\color{darkaccent}Алгоритм: как выбрать способ решения}\par
\vspace{4mm}
{\small\color{softgray}Используйте структуру уравнения до раскрытия громоздких скобок.}
\end{center}
\vspace{5mm}
\centering
\begin{tikzpicture}[
  >=Latex, line width=0.8pt, color=textgray,
  start/.style={ellipse,draw=accent,fill=white,align=center,minimum height=10mm,text width=45mm,font=\small},
  process/.style={rectangle,rounded corners=2mm,draw=accent,align=center,minimum height=12mm,text width=49mm,font=\small},
  test/.style={diamond,aspect=2.0,draw=accent,align=center,inner sep=1pt,text width=34mm,font=\small},
  flow/.style={-{Latex[length=2.2mm]},line width=0.85pt}
]
\node[start] (start) at (0,0) {Перенесите всё влево: $F(x)=0$};
\node[test] (sum) at (0,-2.3) {Сумма квадратов равна $0$?};
\node[process] (system) at (5.7,-2.3) {Решите систему: каждый квадрат равен $0$};
\node[test] (product) at (0,-5.1) {Вся левая часть --- произведение?};
\node[process] (factorzero) at (5.7,-5.1) {Приравняйте каждый множитель к нулю};
\node[process] (factor) at (0,-7.8) {Ищите общий множитель, группировку, формулы сокращённого умножения};
\node[test] (ready) at (0,-10.5) {Удалось получить произведение?};
\node[process] (other) at (0,-13.2) {Нужен другой метод решения};
\node[start] (finish) at (5.7,-13.2) {Найдите все корни и запишите ответ};
\draw[flow] (start) -- (sum);
\draw[flow] (sum) -- node[above,font=\scriptsize]{да} (system);
\draw[flow] (sum) -- node[left,font=\scriptsize]{нет} (product);
\draw[flow] (product) -- node[above,font=\scriptsize]{да} (factorzero);
\draw[flow] (product) -- node[left,font=\scriptsize]{нет} (factor);
\draw[flow] (factor) -- (ready);
\draw[flow] (ready) -| node[near start,above,font=\scriptsize]{да} (factorzero.south);
\draw[flow] (ready) -- node[left,font=\scriptsize]{нет} (other);
\draw[flow] (other) -- (finish);
\draw[flow] (system.east) -- ++(2.2,0) |- (finish.east);
\draw[flow] (factorzero.east) -- ++(1.45,0) |- (finish.east);
\end{tikzpicture}
\par
\clearpage

\topic{Тренировочный блок: Путь А}
\textbf{10 заданий.} Решите уравнения. Оформляйте преобразования последовательно, показывайте вынесенные множители и записывайте ответ по возрастанию корней.

\begin{enumerate}[leftmargin=9mm,itemsep=14pt]
\item \(x^3=x^2-5x+5.\)
\item \(x^3-4x^2-9x+36=0.\)
\item \(x^3-6x^2+5x=0.\)
\item \((x-2)(x+1)(x-5)-(x-4)(x+1)(x-5)=0.\)
\item \(x(x^2-6x+9)=4(x-3).\)
\item \(x(x^2+6x+9)=4(x+3).\)
\item \(x^4=(3x-4)^2.\)
\item \((3x-2)^2=(1-3x)^2.\)
\item \((x^2-3x+2)^2=(x^2-5x-4)^2.\)
\item \((x^2-16)^2+(x^2+3x-4)^2=0.\)
\end{enumerate}

\vfill
\noindent{\bfseries\color{darkaccent}Самопроверка.} Убедитесь, что переносы и раскрытие скобок выполнены с верными знаками. Проверьте вынесенные множители; ни один потенциально нулевой множитель нельзя сокращать. Для системы найдите общие корни, для нулевого произведения --- корни каждого множителя. Запишите все действительные корни по возрастанию.
\clearpage

\topic{Ответы к тренировочному блоку}
\renewcommand{\arraystretch}{1.45}
\begin{tabularx}{\linewidth}{@{}>{\bfseries}p{18mm}X@{}}
\toprule
№ & Корни уравнения (по возрастанию)\\
\midrule
1 & \(1\)\\
2 & \(-3;\ 3;\ 4\)\\
3 & \(0;\ 1;\ 5\)\\
4 & \(-1;\ 5\)\\
5 & \(-1;\ 3;\ 4\)\\
6 & \(-4;\ -3;\ 1\)\\
7 & \(-4;\ 1\)\\
8 & \(\frac12\)\\
9 & \(-3;\ 2-\sqrt5;\ 2+\sqrt5\)\\
10 & \(-4\)\\
\bottomrule
\end{tabularx}

\vspace{3mm}
\textbf{Ключевые промежуточные преобразования для самопроверки}

\begin{enumerate}[leftmargin=9mm,itemsep=2pt]
\item \((x-1)(x^2+5)=0\).
\item \((x-4)(x^2-9)=0\).
\item \(x(x-1)(x-5)=0\).
\item \(2(x+1)(x-5)=0\).
\item \((x-3)(x-4)(x+1)=0\).
\item \((x+3)(x+4)(x-1)=0\).
\item \((x^2-3x+4)(x^2+3x-4)=0\).
\item \((6x-3)(-1)=0\).
\item \((2x+6)(2x^2-8x-2)=0\).
\item \(\begin{cases}x^2-16=0,\\ x^2+3x-4=0.\end{cases}\)
\end{enumerate}

\end{document}
